\origpage{357}{N. Robertson, P.D. Seymour / Journal of Combinatorial Theory, Series B 90 (2004) 325--385}
\hspace*{12pt}$\bullet$ for each $e \in E(\Gamma_i \cap \Sigma_0)$, $\sigma(e) \in E(\Gamma_j \cap \Sigma_0)$ and $\phi_i(e) \le \phi_j(\sigma(e))$,

\hspace*{12pt}$\bullet$ for each circuit $C$ of $F$ and each $v \in \bar{\pi}_C$, $\sigma(v)$ contains $v$.

If we cut $\Sigma_\chi$ along $U(F)$ we obtain two or three fragments; let $\Sigma'$ be the fragment with $\psi(\Sigma') \subseteq \Sigma_0$, where $\psi$ is the associated surjection. Let $\chi'$ be defined as in the proof of 10.2. Since $\chi$ satisfies $\mathbf{S}_3$ and since $\hat{\Sigma}_{\chi'} \cong \hat{\Sigma}_\chi$ and $\Omega_{\chi'}(k) = \Omega_\chi(k)$ ($k = 2, 3$), and $c(\Sigma_{\chi'}) < c(\Sigma_\chi)$, it follows that $\chi'$ is not orientedly bad. For each $i \ge 1$ define $\Gamma'_i, \phi'_i$ as in the proof of 10.2. The sequence $(\Gamma'_i, \phi'_i)$ ($i = 1, 2, \dots$) is similarly oriented, because $\mathscr{S}$ is similarly oriented and if $\Sigma_{\chi'}$ is orientable then so is $\Sigma_\chi$. Thus there exist $j > i \ge 1$ and a linear inflation of $(\Gamma'_i, \phi'_i)$ in $(\Gamma'_j, \phi'_j)$. The result follows as in the proof of 10.2. $\square$

We say that $(\Gamma, \phi, \lambda, \mathscr{T}) \in \mathscr{D}$ is \textit{$n$-flawed in freedom} if $rep(\mathscr{T}) > n$, and $d(r^*_1, r^*_2) > n$ for every two distinct poles $r^*_1, r^*_2$, and there is a circuit $F$ of $\Gamma^*$ with $\frac{1}{2}|E(F)| \le n$ such that $ins(F)$ includes a long side of $\Phi_\chi$ or more than $\frac{1}{2}|E(F)|$ vertices of $\Phi_\chi$. Then $F$ may be chosen to be $\lambda$-linked to $\mathscr{T}$ (by choosing $F$ with $ins(F)$ maximal), and in this case we call it a \textit{freedom flaw} for $(\Gamma, \phi, \lambda, \mathscr{T})$. By adapting the proofs of 10.2, 10.3, 10.4 in the natural way, using $\mathbf{S}_4$ in place of $\mathbf{S}_1$ and $\mathbf{S}_3$, we obtain

\textbf{10.5.} \textit{Let $\chi$ satisfy $\mathbf{S}_4$ \fix{and $\mathbf{S}_3$}, and let $n \ge 0$ be an integer. Then there is a well-behaved set $\mathscr{C}_3(n)$ of $\chi$-places with the following property. Let $(\Gamma, \phi, \lambda, \mathscr{T}) \in \mathscr{D}$ be $n$-flawed in freedom. Then there is a rooted location $\mathscr{L}$ which $(n + 1)$-isolates $\mathscr{T}$ and for which $(\Gamma, \phi, \mathscr{L}) \in \mathscr{C}_3(n)$.}\fixnote{fix:10.5}{Unlike the flaws of 10.3 and 10.4, a freedom flaw comes with no lower bound on the distance between the cuffs. So the circuit $F$ may enclose a long side on one cuff and meet a second cuff in one vertex. Cutting along $F$ then joins $U(F)$ and the second cuff into one cuff, and the outer piece has fewer cuffs than $\Sigma_\chi$. A refinement of $\chi$ has as many cuffs as $\Sigma_\chi$, so $\mathbf{S}_4$ does not apply to that piece, and $\mathbf{S}_3$ is needed. For example, let $n=3$, and let $\Sigma_\chi$ be an annulus with cuffs $\Theta_1,\Theta_2$. Let $\Theta_1$ have two frame vertices $a,b$, with one long side $L$ and one short side, and let $\Theta_2$ have two frame vertices $c,d$, with two short sides. Let $\Gamma$ consist of a large triangulated grid with a hole bounded by four vertices $c,d,u,w$, where $\Theta_2$ runs along $c,d$, and of a wide triangulated disc around $\Theta_1$, attached to the grid at $u$, $w$ and $c$. Then $d(r_1^*,r_2^*)=4>3$, and the freedom flaw $F$ with $ins(F)$ maximal has length $6$. It passes through $c\in\Theta_2$ and encloses $L\subseteq\Theta_1$, so cutting along $F$ leaves one cuff instead of two. The only use of 10.5 is in the proof of 4.1, where $\mathbf{S}_3$ holds, so nothing later is lost.}

\fix{We give the proof, by adapting the proofs of 10.2, 10.3 and 10.4; we first choose the circuit $F$ with care. Let $(\Gamma, \phi, \lambda, \mathscr{T}) \in \mathscr{D}$ be $n$-flawed in freedom. Call a circuit $F$ of $\Gamma^*$ \textit{eligible} if $\frac{1}{2}|E(F)| \le n$ and $ins(F)$ includes a long side of $\Phi_\chi$ or more than $\frac{1}{2}|E(F)|$ vertices of $\Phi_\chi$. (Since $\frac{1}{2}|E(F)| \le n < rep(\mathscr{T})$, $ins(F)$ is defined.) Choose an eligible $F$ with $|E(F)|$ minimum, and subject to this with $|V(\Gamma^*) \cap ins(F)|$ maximum. By 8.7 there is a circuit $F'$ which is $\lambda$-linked to $\mathscr{T}$, with $|E(F')| \le |E(F)|$ and $ins(F) \subseteq ins(F')$; then $F'$ is eligible, and it satisfies the same two extremal conditions, so we may assume that $F$ is $\lambda$-linked to $\mathscr{T}$.

(a) \textit{Let $Q$ be a path of $\Gamma^*$ with ends $a, b \in V(F)$ and with no point in the interior of $ins(F)$, and let $A, B$ be the two paths of $F$ between $a$ and $b$. If $|E(Q)| \le |E(A)|, |E(B)|$, then $V(Q) \subseteq ins(F)$ and $|E(Q)|$ equals one of $|E(A)|, |E(B)|$.} For one of the circuits $Q \cup A$, $Q \cup B$, say $X$, has $ins(X)$ meeting the interior of $ins(F)$, and then $ins(F) \subseteq ins(X)$, since $U(X)$ is disjoint from the interior of $ins(F)$. So $X$ is eligible, and $|E(X)| \le |E(F)|$. If $Q$ had a vertex not in $ins(F)$, or were shorter than both $A$ and $B$, then $X$ would contradict the choice of $F$.

(b) \textit{Exactly one pole $r^*$ is in $ins(F)$, and the closed disc $\bar{r}'$ of every other pole $r'^*$ meets $ins(F)$ in at most one point.} If two poles were in $ins(F)$, their distance would be at most $\frac{1}{2}|E(F)| \le n$. Let $r'^*$ be a pole not in $ins(F)$. Then $r'^* \notin V(F)$, so $r'$ is disjoint from $U(F)$, and since $r'^* \notin ins(F)$ it follows that $\bar{r}' \cap ins(F) \subseteq U(F) \cap V(\Gamma)$; and two vertices $u, v$ of $\Gamma$ in this set give a path $Q$ with edges $ur'^*$, $r'^*v$ that contradicts (a), since $|E(A)|, |E(B)| \ge 2$ as $\Gamma^*$ is bipartite. If no pole were in $ins(F)$, then by this argument $ins(F)$ would include no long side, and the vertices of $\Phi_\chi$ in $ins(F)$ would lie in $U(F) \cap V(\Gamma)$, which has $\frac{1}{2}|E(F)|$ members; so $F$ would not be eligible.

(c) Let $\Theta = \bar{r} \setminus r$, and let $K$ be the drawing in $\Sigma_\chi$ with $U(K) = U(F) \cap \Sigma_\chi$ and $V(K) = U(F) \cap V(\Gamma)$. Then $K$ is a feature, and $U(K) \cap U(\Gamma) = V(K)$. Since $\Gamma^*$ is bipartite, the vertices of $F$ are alternately vertices of $\Gamma$ and poles, and so $|V(K)| \le \frac{1}{2}|E(F)|$. Moreover $K$ has at most $|V(K)|$ edges, since $U(K)$ is either the $O$-arc $U(F)$ (if $r^* \notin V(F)$) or a line with both ends in $V(K)$ (if $r^* \in V(F)$, since then the two edges of $F$ at $r^*$ lie in $\bar{r}$ and $r^*$ is the only pole in $V(F)$, by (b)). Hence the size of $K$ is at most $2|V(K)| \le |E(F)| \le 2n$. If $r^* \notin V(F)$, let $Z = ins(F)$; then $\bar{r} \subseteq Z$, since $r$ is disjoint from $U(F)$ and $r^* \in ins(F)$, and $bd(Z) = U(K)$. If $r^* \in V(F)$, the two edges of $F$ at $r^*$ form a line which cuts $\bar{r}$ into two closed discs, one in $ins(F)$, and we let $Z$ be the union of $ins(F)$ and the other one. Then $Z$ is a closed disc including $\bar{r}$, and $bd(Z) = U(K) \cup A$ where $A$ is an arc of $\Theta$. Since every other pole is outside $ins(F)$, it is outside $Z$, and by (b) its closed disc meets $Z$ in at most one point.

(d) Cut $\Sigma_\chi$ along $U(K)$, and let $\Sigma'$ be the fragment with $\psi(\Sigma') = \Sigma_\chi \setminus (Z \setminus bd(Z))$. Capping $\Sigma'$ gives a surface homeomorphic to $\hat{\Sigma}_\chi$, since $Z$ is a disc; and if $\Sigma'$ is orientable then so is $\Sigma_\chi$. Let $\chi'$ be defined from $\Sigma'$ as in the proof of 10.2. Since $\Theta \subseteq Z$, the cuffs of $\Sigma'$ are the cuffs of $\Sigma_\chi$ disjoint from $Z$, and one more cuff, which lies over $bd(Z)$ and over the other cuffs that meet $Z$. So if some cuff other than $\Theta$ meets $Z$, then $c(\Sigma') < c(\Sigma_\chi)$, and $\chi'$ is not orientedly bad, by $\mathbf{S}_3$. Otherwise $c(\Sigma') = c(\Sigma_\chi)$, so $\Sigma' \cong \Sigma_\chi$ by the classification of surfaces, and $\psi$ is one-to-one on $\Sigma'$, since near each point of $bd(Z)$ the complement of $Z$ is connected. We show that then $\chi'$ is a refinement of $\chi$, with the embedding $f$ that maps each long side $R$ of $\Phi_{\chi'}$ to the long side of $\Phi_\chi$ including $\psi(R)$; so $\chi'$ is not orientedly bad, by $\mathbf{S}_4$. Let $I = (Z \setminus bd(Z)) \cup U(K)$. Since $F$ is eligible and $ins(F) \cap \Sigma_\chi \subseteq I$, either $I$ includes a long side, or $|V(\Phi_\chi) \cap I| > \frac{1}{2}|E(F)| \ge |V(K)|$. These long sides and vertices are in $\Theta$, since the other cuffs miss $Z$; and so some vertex $w$ of $\Phi_\chi$ in $\Theta$ is not in $\Theta \setminus I$, which is an open arc of $\Theta$ or empty. Hence for each long side $S$ in $\Theta$, $S \setminus I$ is connected, since $S$ and $\Theta \setminus I$ are arcs of $\Theta$ that miss $w$; and the long sides in the other cuffs do not meet $I$. So $f$ is one-to-one, and the first two conditions of a refinement hold. If $f$ is onto, then $I$ includes no long side, so $|V(\Phi_\chi) \cap I| > |V(K)|$, and since $\psi$ is one-to-one on $\Sigma'$, $|V(\Phi_{\chi'})| \le |V(\Phi_\chi)| - |V(\Phi_\chi) \cap I| + |V(K)| < |V(\Phi_\chi)|$. A frame has as many sides as vertices, since $k \ge 1$ vertices cut a cuff into $k$ arcs; and $\Phi_{\chi'}$ and $\Phi_\chi$ have the same number of long sides. Hence $\Phi_{\chi'}$ has fewer short sides than $\Phi_\chi$, and the third condition holds.

(e) Now we adapt the proofs of 10.2, 10.3 and 10.4, with the following changes. In place of the finitely many pairs $(F_i, \pi_i)$ of the proof of 10.3 we use 5.2: there are features $K_1, \dots, K_m$ in $\Phi_\chi$, each of size at most $2n$, such that every feature in $\Phi_\chi$ of size at most $2n$ is equivalent to one of them. For $1 \le c \le m$, cutting $\Sigma_\chi$ along $U(K_c)$ gives finitely many fragments. In place of the sets $\mathscr{C}(F_i, \pi_i)$ of 10.2 we use the sets $\mathscr{C}(K_c, \Sigma')$ defined below, one for each $c$ and each fragment $\Sigma'$ obtained by cutting along $U(K_c)$; and as in the proof of 10.3 we let $\mathscr{C}_3(n)$ be the union of the sets of all images of members of these sets under $\Phi_\chi$-preserving homeomorphisms. By 9.4, $\mathscr{C}_3(n)$ is well-behaved, since there are only finitely many sets $\mathscr{C}(K_c, \Sigma')$, provided that each of them is well-behaved, which we show below. Now let $(\Gamma, \phi, \lambda, \mathscr{T}) \in \mathscr{D}$ be $n$-flawed in freedom, and let $F, K, Z, \Sigma'$ be as in (a)--(d). By (c), $K$ is equivalent to some $K_c$, by a $\Phi_\chi$-preserving homeomorphism $\alpha$ with $\alpha(K) = K_c$. As in the proof of 10.3, the image of $(\Gamma, \phi, \lambda, \mathscr{T})$ under $\alpha$ is in $\mathscr{D}$ and is $n$-flawed in freedom, and $\alpha(F)$ and $\alpha(Z)$ have for it the properties of $F$ and $Z$ above; so, replacing $(\Gamma, \phi, \lambda, \mathscr{T})$ by its image and the rooted location found by its image under $\alpha^{-1}$ at the end, we may assume that $K = K_c$.

Fix $c$ and a fragment $\Sigma'$ obtained by cutting $\Sigma_\chi$ along $U(K_c)$, with surjection $\psi$, and let $\pi$ be a march with $\bar{\pi} = V(K_c)$. Let $\mathscr{D}''$ be the set of all $(\Gamma, \phi, \lambda, \mathscr{T}) \in \mathscr{D}$ that are $n$-flawed in freedom and have a circuit $F$ as above with $K = K_c$ and with $\Sigma'$ as the fragment of (d). For each of them, let $A$ be the rooted hypergraph with $A^- = \Gamma \cap ins(F)$ and $\pi(A) = \pi$, and let $\mathscr{L} = \{A\}$; here $\bar{\pi} = V(K) = U(F) \cap V(\Gamma)$. Then $\mathscr{L}$ $(n+1)$-isolates $\mathscr{T}$, by 8.8, since $F$ is $\lambda$-linked to $\mathscr{T}$, $|\bar{\pi}| \le \frac{1}{2}|E(F)| \le n$ and $n + 1 \le rep(\mathscr{T}) \le ord(\mathscr{T})$. Let $\mathscr{C}(K_c, \Sigma')$ be the set of all such $(\Gamma, \phi, \mathscr{L})$. Let $\Sigma_2 = \psi(\Sigma')$. As in the proof of 10.2, by 9.3 it suffices to show that for every countable sequence of members of $\mathscr{C}(K_c, \Sigma')$ there exist $j > i \ge 1$ and an inflation $\sigma$ of $\Gamma_i \cap \Sigma_2$ in $\Gamma_j$ satisfying the three conditions listed in the proof of 10.2. In place of the colour scheme of the proof of 10.2, which is not bad by $\mathbf{S}_1$, we use the scheme $\chi'$ of (d), which is not orientedly bad by $\mathbf{S}_3$ or by $\mathbf{S}_4$. Define $(\Gamma'_i, \phi'_i)$ as in the proof of 10.2; these are $\chi'$-coloured paintings. As in the proof of 10.4, the sequence $(\Gamma'_i, \phi'_i)$ $(i = 1, 2, \dots)$ is similarly oriented, because $\mathscr{S}$ is similarly oriented and, by (d), if $\Sigma'$ is orientable then so is $\Sigma_\chi$. Hence there exist $j > i \ge 1$ and a linear inflation $\sigma'$ of $(\Gamma'_i, \phi'_i)$ in $(\Gamma'_j, \phi'_j)$. In place of $\sigma(v) = \psi(\sigma'(v'))$ where $\psi^{-1}(v) = \{v'\}$, as in the proof of 10.2, we let $\sigma(v)$ be the image under $\psi$ of the union of the drawings $\sigma'(v')$ over all $v' \in \Sigma'$ with $\psi(v') = v$, and $\sigma(e) = \psi(\sigma'(\psi^{-1}(e) \cap \Sigma'))$ for each edge $e$; there may be two points $v'$ over a vertex $v$ of $K_c$ in a cuff other than $\Theta$, when $\mathbf{S}_3$ is used in (d). Then $\sigma$ is an inflation of $\Gamma_i \cap \Sigma_2$ in $\Gamma_j$ satisfying the three conditions, by the argument given in detail at the end of the proof of 5.4, where the assumption that $V(K) \cap V(\Gamma_i) = V(K) \cap V(\Gamma_j)$ holds here because $V(K_c) \subseteq V(\Gamma_i) \cap V(\Gamma_j)$. In particular each $v \in \bar{\pi}$ is in $V(\sigma(v))$, since every point over $v$ is a vertex of $\Phi_{\chi'}$ and $\sigma'$ respects $\Phi_{\chi'}$. So $\mathscr{C}(K_c, \Sigma')$ is well-behaved, and we have shown that $(\Gamma, \phi, \mathscr{L}) \in \mathscr{C}_3(n)$ for every $(\Gamma, \phi, \lambda, \mathscr{T})$ that is $n$-flawed in freedom. This proves 10.5.}\fixnote{fix:10.5-proof-added}{The paper omits the proof of 10.5, saying only that it is obtained by adapting the proofs of 10.2, 10.3 and 10.4. The argument above is the one used in the formalization.}

\sectionhead{11. Border label flaws}

So far we have examined flaws in our paintings $(\Gamma, \phi)$ due to some kind of lack of generality in $\Gamma$---representativeness, distance and freedom flaws. There are two other flaws we must investigate, both concerned with a lack of generality in $\phi$, and it is to deal with these that $\mathbf{S}_2$ and the full strength of $\mathbf{S}_4$ are needed. In this section we examine the situation when for some long side $S$, the values of $\phi(e)$ over edges $e$ bordering $S$ fail to be sufficiently general.

Let $\chi, \mathscr{S}, \mathscr{D}$ be as in Section 10. Let $(\Gamma, \phi, \lambda, \mathscr{T}) \in \mathscr{D}$, and let $r^*$ be a pole. By a \textit{bite at $r^*$} we mean a circuit $F$ of $\Gamma^*$ with $\frac{1}{2}|E(F)| < rep(\mathscr{T})$, such that $r^* \in V(F)$. If $e_1, e_2$ are edges of $\Gamma$ bordering $\bar{r}\setminus r$, we define $l(e_1, e_2)$ to be the minimum of $\frac{1}{2}|E(F)|$ over all bites $F$ at $r^*$ with $e_1, e_2 \subseteq ins(F)$, if there is such a bite, and otherwise $l(e_1, e_2) = rep(\mathscr{T})$.

Let $S$ be a long side of $\Phi_\chi$, let $m \ge 1$, let $\omega_1, \dots, \omega_m$ be a sequence of elements of $\Omega_\chi(S)$, and let $n \ge 4$. We say that $(\Gamma, \phi, \lambda, \mathscr{T}) \in \mathscr{D}$ is $(n, (\omega_1, \dots, \omega_m))$-\textit{flawed on $S$} if

\hspace*{12pt}$\bullet$ $rep(\mathscr{T}) \ge 2(m + 1)n + 8$, and $dist(\Gamma) \delm{\ge}\fix{>} \frac{1}{2}(m + 1)n + 1$,\fixnote{fix:11-flawed-on-S}{Being $N$-flawed in distance asks $rep(\mathscr{T}) > N$ and $dist(\Gamma) > \frac12 N + 1$ (page 355). As printed, $dist(\Gamma) = \frac{1}{2}(m+1)n + 1$ is allowed here, and then ``not $(m+1)n$-flawed in distance'' holds for a trivial reason. It then gives no lower bound on the distance between poles, which the proof of 11.2 needs. Claim (1) in the proof of 4.1, the only place that produces this flaw, gives the strict bound.}

\hspace*{12pt}$\bullet$ $(\Gamma, \phi, \lambda, \mathscr{T})$ is not $(m + 1)n$-flawed in distance,

\hspace*{12pt}$\bullet$ $(\Gamma, \phi, \lambda, \mathscr{T})$ is not $(m + 1)n$-flawed in freedom, and
